%! Author = Saeid Yazdani
%! Date = 7/7/2024

\section{Effects of Current Limiting Circuits}
\label{sec:current-limiting-effects}

As stated before, in a typical \ac{ATE} setup, power supplies provide current to
both the DUT and its decoupling capacitors as well as parasitic capacitances.
The current going to the capacitors can significantly be based on the total
capacitance and power-up time set by the test developer.
As an example, powering up a \SI{3.3}{\volt} \ac{DUT} in \SI{100}{\micro\second}
with a total circuit capacitance of \SI{20}{\micro\farad} can need up to
\SI{1}{\ampere} according to \autoref{eq:current-during-capacitor-charge}.

\begin{equation}
    \label{eq:current-during-capacitor-charge}
    I = C \frac{dV}{dt}
\end{equation}

Almost all modern \ac{DPS} provide a programmable current limit option,
typically in \SI{10}{\milli\ampere} to \SI{50}{\milli\ampere} steps.
During test development and execution, it is a common practice to protect
\acp{DUT} from uncontrolled events by setting limits for the maximum current
a \ac{DPS} is allowed to source into a \ac{DUT}.

If a brief surge in current triggers a preset limit, the circuit will turn off
the limit once the surge stops, and the DPS will go back to providing the set
voltage with a delay.
Depending on the voltage level before the occurrence of a current surge event,
the actual voltage might need to go higher than the set voltage to make sure the
problem that caused the high current is settled.
Therefore the DUT might experience a voltage higher than the set
voltage for a short period.

If the test developer sets tight current limits, especially during the
\emph{TEST START} where all the capacitors need to be charged, there is a high
chance for misbehavior of the voltage rails due to the engagement of the current
limit circuitry.
This explains the presence of glitches
on \autoref{fig:villach-supply-analysis-1v3-test-start-glitch} and other
similar measurements.

\subsection{Mitigating Current-Limit Induced Glitching}
\label{subsec:current-limiting-effects}

As stated before, current depends on capacitance and how fast the voltage
changes, so the highest risk of engaging the current limit and overshooting
happens during power-up.
Test developers need to check power-up steps and the current limit.
Changing voltages by small amounts are less likely to be affected by the
capacitance and current limit because the voltage change is smaller.

The following points may help with reducing the overshoot and undershoot
triggered by the current limiting feature of a \ac{DPS}:

\begin{itemize}[topsep=8pt,itemsep=4pt,partopsep=4pt, parsep=4pt]
    \item \textbf{Relaxing the current limit:}
    Depending on internal circuitry and control loop of \ac{DPS} modules, test
    developers might need to increase current limits by at
    least \SI{100}{\milli\ampere} per microfarad to reduce the chance of
    high voltage glitches due to immature activation of current limit circuits.

    \item \textbf{Stepwise increase to target voltage:}
    Many modern \ac{DPS} allow for voltage step-ups with fine control over time
    taken by each step.
    If the \ac{DUT} power requirements allow, test developers can program the
    \ac{DPS} to reach target voltage in multiple steps, instead of one.

\end{itemize}












